\section{The Indigenous Cosmologies in Comparison}

This part has held eighteen cosmologies drawn from the Americas, Asia and the Pacific, Africa and the Arctic, and Europe and Persia. They differ widely in age, scale, and degree of documentation. Set side by side, the same small set of creation moves recurs across them, the same moves already seen in the ancient and the esoteric layers of this collection. This section lays out the shared moves and groups the eighteen by the move that dominates each.

\subsection{The Shared Moves}

\textbf{The primordial water or void.} Many begin from a wide undivided state that is not nothing. It is a marsh of water, a deep darkness, a fused mass, or a formless quiet. The world starts from a full but undifferentiated source, not from emptiness.

\textbf{Emergence through worlds or the earth-diver.} In one family of accounts life climbs upward through stacked worlds toward the present one. In another a being descends or dives to bring up earth and spread it over the first water to make solid land.

\textbf{The separation of sky and earth.} A joined or fused whole is pulled into two and held apart, and the gap between them becomes the room of the world.

\textbf{The high god and the trickster or culture-bringer.} A remote supreme power stands behind creation, while a nearer figure does the shaping work, brings light or fire, founds the order, or upsets it. The two roles are often split between distinct beings.

\textbf{The layered cosmos.} The finished world is built in levels, an upper sky, a middle ground, and a lower or sea-floor underworld, with the dead and the powers assigned to their tiers.

\subsection{The Eighteen by Dominant Motif}

\begin{longtable}{L{0.34\linewidth}
L{0.56\linewidth}}
\toprule
\textbf{Dominant motif} & \textbf{Traditions} \\
\midrule
\textbf{Separation of sky and earth} & Polynesian, Korean, Chinese-influenced layers \\
\textbf{Emergence through stacked worlds} & Navajo, Zuni, Hopi \\
\textbf{Earth-diver, land from the water} & Yoruba \\
\textbf{Primordial water or deep darkness} & Polynesian, Inuit, Maya \\
\textbf{Dual powers in conflict} & Zoroastrian, Aztec, Slavic, Baltic \\
\textbf{Shaping of an unformed land by ancestral beings} & Aboriginal, Andean \\
\textbf{Birth or fashioning from a divine maker} & Tamil, Yoruba, Finnish, Celtic \\
\bottomrule
\end{longtable}

The groupings are not exclusive. Several traditions carry more than one move and appear under more than one heading. The table marks only the move that stands out most in each.

\subsection{The Go-Out-and-Return}

Under all eighteen runs the deep theme this collection traces everywhere. The \textbf{going-out} and the \textbf{return}. One source goes out into the many. The water or darkness or fused mass differentiates into sky, earth, land, powers, and people. That going-out is present in every account.

The return takes different forms. In some it is cosmic, a world that ends and renews. In some it is ritual, the spirit-journey that sets the broken order right and is run again whenever it is needed. In some it is personal, the soul aligning with its destiny or turning inward to find the source already present in the land and the body. The eighteen agree on the shape. One source, going out into the many, with a return granted to the cosmos, to the rite, or to the soul.
